\documentclass[12pt, a4paper]{article}
\usepackage[utf8]{inputenc}
\usepackage[T2A]{fontenc}
\usepackage[normalem]{ulem}
\usepackage{ragged2e}
\usepackage[style=russian]{csquotes}
\usepackage[english,russian]{babel}
\usepackage{graphicx,wrapfig}
\usepackage{hyperref}
\usepackage{amsmath}
\usepackage{tikz}
\usepackage{caption}
\captionsetup[table]{justification=raggedright,singlelinecheck=off}
\captionsetup[lstlisting]{justification=raggedright,singlelinecheck=off}
\captionsetup{labelsep=endash}
\captionsetup[figure]{name={Рисунок}}
\usepackage{listings}
\usepackage{indentfirst, setspace}
\usepackage{tabularx, multirow}
\linespread{1.3}
\lstset{
  language=C++, % Язык программирования (можно выбрать другой)
  basicstyle=\ttfamily\small, % Основной стиль шрифта
  keywordstyle=\color{blue}, % Стиль ключевых слов
  breaklines=true,
  commentstyle=\color{green}, % Стиль комментариев
  stringstyle=\color{red}, % Стиль строк
  extendedchars=\true,
  showstringspaces=false,
}
\usepackage[left=20mm, top=15mm, right=15mm, bottom=20mm]{geometry}

\begin{document}
\begin{titlepage}
     \fontsize{12}{12}\selectfont

  {\centering

   \begin{bf}

    \begin{wrapfigure}{l}{19mm}
        
\includegraphics[width=22mm]{emblem.jpeg}
    \end{wrapfigure}


    \noindent Министерство науки и высшего образования Российской Федерации

    \noindent Федеральное государственное бюджетное образовательное учреждение высшего образования

    \noindent \enquote{Московский государственный технический унверситет

     \noindent имени Н.Э. Баумана

     \noindent (национальный исследовательский университет)}

    \noindent (МГТУ им. Н.Э. Баумана)

   \end{bf}
  }

  \vspace{0.4cm}

  {\setstretch{0.1}
   \noindent\rule{\textwidth}{1mm}
   \noindent\rule{\textwidth}{0.5mm}

  }

  \fontsize{14}{21}\selectfont

  \noindent\begin{tabularx}{\textwidth}{l >{\centering\arraybackslash}X}
   ФАКУЛЬТЕТ & \flqq Фундаментальные Науки\frqq \\ \cline{2-2}

   КАФЕДРА & ФН-12 \flqq Математическое моделирование\frqq \\ \cline{2-2}
  \end{tabularx}


  \vspace{1cm}


  \begin{center}
   \begin{bf}

    \fontsize{24}{36}\selectfont
    ОТЧЕТ

    \fontsize{20}{30}\selectfont
    ПО ЛАБОРАТОРНОЙ РАБОТЕ НА ТЕМУ:

    Редакционные расстояния

   \end{bf}
  \end{center}

  \fontsize{14}{21}\selectfont
  \vspace{5cm}


  \noindent\begin{tabularx}{\textwidth}{ X >{\centering}p{4cm} p{1cm} c }
   Студент: & & & Подобедов В. В. \\ \cline{2-2} \cline{4-4}
   & \fontsize{10}{15}\selectfont дата, подпись & & \fontsize{10}{15}\selectfont Ф.И.О. \\
   Преподаватель: & & & Волкова Л. Л.\\ \cline{2-2} \cline{4-4}
   & \fontsize{10}{15}\selectfont дата, подпись & & \fontsize{10}{15}\selectfont Ф.И.О.
  \end{tabularx}

  \vspace{\fill}

  \begin{center}
   \it{Москва}, 2023
  \end{center}

  \thispagestyle{empty}
\end{titlepage}


\renewcommand{\contentsname}{Содержание}
\setcounter{page}{2}
\tableofcontents

\newpage
\begin{center}
   \section*{Введение} 
\addcontentsline{toc}{section}{Введение}
\end{center}


\vspace{5mm}
\raggedright
\justifying
\textbf{Цель лабораторной работы}: провести сравнительный анализ алгоритмов поиска редакционных расстояний.


Для достижения поставленной цели требуется решить следующие \textbf{задачи}.
\begin{enumerate}
\item Описать расстояния Левенштейна и Дамерау-Левенштейна.
\item Разработать алгоритмы вычисления редакционных расстояний, согласно варианту.
\item Реализовать разработанные алгоритмы.
\item Выполнить тестирование реализации разработанных алгоритмов.
\item Разработать программное обеспечение с двумя режимами работы --- режимом расчёта расстояний для введённых пользователем двух строк произвольной длины и режимом массированного замера процессорного времени для автоматически сгенерированных строк заданной длины.
\item Провести сравнение временной эффективности реализации разработанных алгоритмов, проведя замеры процессорного времени их выполнения в зависимости от линейного размера строк одинаковой длины.
\item Выполнить оценку емкостной эффективности разработанных алгоритмов в зависимости от линейных размеров строк (для рекурсивных алгоритмов дать оценку пиковой затрачиваемой памяти) Провести сравнение временной эффективности реализации разработанных алгоритмов, проведя замеры памяти, затрачиваемой реализацией разработанных алгоритмов в зависимости от линейных размеров строк (для рекурсивных алгоритмов замерить пиковую затрачиваемую память).
\end{enumerate}

Согласно варианту, требуется разработать следующие три алгоритма.
\begin{enumerate}
\item Алгоритм вычисления расстояния Левенштейна нерекурсивным методом, используя две строки.
\item Алгоритм вычисления расстояния Дамерау-Левенштейна нерекурсивным методом, пользуясь матрицей и в конце программы вывести её содержимое.
\item Алгоритм вычисления расстояния Дамерау-Левенштейна рекурсивным методом без кэша.
\end{enumerate}


\newpage
\begin{center}
   \section{Аналитическая часть} 
\end{center}

Расстояние между строками a и b можно найти при помощи некоторых редакционных расстояний. Рассмотрим их.

Расстояние Левенштейна --- это метрика cходства между двумя строковыми последовательностями \cite{mpl}. Оно может быть рассчитано по формуле ~\ref{eq:Lev}:\\
\begin{equation} \label{eq:Lev}
d_{a, b}(i,j) =
\begin{cases}
    0 & i = 0, j = 0
\\   i & j = 0, i > 0
\\   j & i = 0, j > 0
\\ \min \begin{cases}
    d_{a, b}(i, j - 1) + 1
\\  d_{a, b}(i - 1, j) + 1
\\  d_{a, b}(i - 1, j - 1) + m(a[i], b[j])
\end{cases}  & j > 0, i > 0  
\end{cases}\\
\end{equation}

Нетрудно понять, что первые три строчки формулы заполнят первый столбец и первую строку матрицы. Для всех остальных ячеек матрицы используется четветрая формула. Стоит уделить внимание оператору ``m(a[i], b[j])''. Он сравнивает символы в слове, и, если они одинаковые, он дает 0, иначе --- 1.


Расстояние Дамерау-Левенштейна вычисляется по формуле ~\ref{eq:Levv}:
\begin{equation} \label{eq:Levv}
d_{a, b}(i, j) = 
\begin{cases} 
\max(i, j) & \mbox{если } i \cdot j = 0;
\\ \min \begin{cases} 
d_{a, b}(i - 1, j) + 1
\\d_{a, b}(i, j - 1) + 1
\\d_{a, b}(i - 1, j - 1 ) + 1_{(a_i = b_j)}
\\d_{a, b}(i - 2, j - 2) + 1
\end{cases} & \mbox{если } i,j > 1, a_i = b_{j - 1}, a_{i - 1} = b_j;
\\ \min \begin{cases} 
d_{a, b}(i - 1, j) + 1
\\d_{a, b}(i, j - 1) + 1
\\d_{a, b}(i - 1, j - 1) + 1_{(a_i = b_j)}
\end{cases} & \mbox{иначе} 
\end{cases}\\
\end{equation}

\justifying
Значительным отличием в нахождении расстояния Дамерау-Левенштейна от расстояния Левенштейна является транспозиция соседних символов. 

Говоря об эффективности методов вычисления расстояния слов, хочется отметить рекурсивный способ Дамерау-Левенштейна с кэшем. Поскольку после каждого подсчета порождается еще три вычисления, рекурсия с кэшем может помочь сохранить процессорное время. Поэтому вместо $3^{k*l}$ раз, где k и l длина слов соответсвенно, нужно вычислить $k * l$ раз.
Что касается других способов решения задачи, можно использовать нерекурсивный метод с заполнением матрицы размера $(k+1)*(l+1)$. Однако можно заметить, что для вычисления расстояния Левенштейна используются только числа предыдущей и вычисляемой строки.  Поэтому в каждый момент вычисления можно хранить две строки матрицы.

\newpage
\begin{center}
    \section{Конструкторская часть}
\end{center}
Cхема алгоритма поиска расстояния Левенштейна нерекурсивным способом с хранением двук строк матрицы представлена на рис.~\ref{img:algLevNR}.
\begin{figure}[h!]
    \center{\includegraphics[width=5cm, height=22cm]{block_1}}
    \caption{Cхема алгоритма поиска расстояния Левенштейна нерекурсивным способом}
    \label{img:algLevNR}
\end{figure}

\newpage
Cхема алгоритма поиска расстояния Дамерау-Левенштейна нерекурсивным способом с хранением матрицы представлена на рис.~\ref{img:algDLevNR}. \\
\begin{figure}[h!]
    \center{\includegraphics[width=5cm, height=22cm]{block_2}}
    \caption{Cхема алгоритма поиска расстояния Дамерау-Левенштейна нерекурсивным способом с хранениемматрицы}
    \label{img:algDLevNR}
\end{figure}

\newpage
Cхема алгоритма поиска расстояния Дамерау-Левенштейна рекурсивным способом с хранением 3 строк матрицы представлена на рис.~\ref{img:algDLevNR2}. \\
\begin{figure}[h!]
    \center{\includegraphics[width=14cm, height=22cm]{block_3}}
    \caption{Cхема алгоритма поиска расстояния Дамерау-Левенштейна рекурсивным способом с хранением 3 строк матрицы}
    \label{img:algDLevNR2}
\end{figure}

\newpage
\begin{center}
    \section{Технологическая часть}
\end{center}

%\subsection{Выбор средств реализации и языка программирования}
Для реализации алгоритмов был выбран язык программирования С++. Кроме стандартной библиотеки <iostream>, которая предоставляет возможность пользоваься срадствами ввода-вывода для стандартной консоли. Библиотека <algorithm> подключена для использования функции min(). Также для генерации различных слов от времени была взаимствована библиотека <ctime>.
\par Функции, разработанные в ходе написания лабораторной работы. 
\begin{enumerate}
\item levenstein --- функция, которая нерекурсивно вычисляет расстояние двух слов, используя две строки по алгоритму Левенштейна.
    \item DamerauLevenshteinDistance --- функция, которая нерекурсивно вычисляет расстояние двух строк, используя для этого матрицу по алгоритму Дамерау-Левенштейна.
    \item damerauLevenshteinDistance --- функция, которая рекурсивно вычисляет расстояние двух слов по алгоритму Дамерау-Левенштейна с проверкой на кэш в матрице. 
\end{enumerate}

%\subsection{2222}
В программе реализовано переключение между двумя режимами: пользовательским и автоматическим. Первый способ дает возможность человеку самому ввести необходимые слова и узнать их расстояние. Второй способ дает возможность отследить ресурсоёмкость каждого из трех алгоритмов.
\par Далее предствлены листинги~\ref{lst:label1}, \ref{lst:label2}, \ref{lst:label3} реализаций трех алгоритмов.


\begin{lstlisting}[caption={Расстояние Левенштейна нерекурсивный метод, используя две строки матрицы},label={lst:label1}]
#include <iostream>
#include <string>
#include <vector>
#include <algorithm>
#include <ctime>

using namespace std;

string RandomWord(int length) {
	static const char alphabet[] = "абвгдеёжзиклмнопрстуфхцчшщъыьэюя";
	string result;

	for (int i = 0; i < length; ++i) {
		int randomIndex = rand() % (sizeof(alphabet) - 1);
		result += alphabet[randomIndex];
	}

	return result;
}

int levenstein(string str_1, string str_2) {
	int n = str_1.length();
	int m = str_2.length();

	vector<int> current_row(n + 1);
	for (int i = 1; i <= n; i++) {
		current_row[i] = i;
	}

	for (int i = 1; i <= m; i++) {
		vector<int> previous_row = current_row;
		current_row[0] = i;
		for (int j = 1; j <= n; j++) {
			current_row[j] = min(previous_row[j] + 1, min(current_row[j - 1] + 1, 
				(str_1[j - 1] != str_2[i - 1]) ? previous_row[j - 1] + 1 : previous_row[j - 1]));
		}
	}
	return current_row[n];
}

int main() {
    setlocale(0, "");
	string str1, str2;
	srand(time(0));
	int a, k, num;
	double medtime = 0;
	cout << "Ввести данные вручную (1) или сгенерировать (2) - ";
	cin >> a;
	if (a == 1) {
		cout << "Введите первое слово: ";
		cin >> str1;
		cout << "Введите второе слово: ";
		cin >> str2;
		clock_t start_time = clock();
		cout << "Расстояние левенштейна: " << levenstein(str1, str2) << endl;
		clock_t end_time = clock();
		double cpu_time = (double(end_time - start_time)) / CLOCKS_PER_SEC;
		cout << "Процессорное время: " << cpu_time << endl;
		cout << endl;
	}
	else {
		cout << "Введите длину слова: ";
		cin >> k;
		cout << "Количество слов: ";
		cin >> num;
		for (int i = 1; i <= num; i++) {
			cout << i << endl;
			str1 = RandomWord(k);
			str2 = RandomWord(k);
			cout << str1 << endl << str2 << endl;
			clock_t start_time = clock();
			cout << "Расстояние левенштейна: " << levenstein(str1, str2) << endl;
			clock_t end_time = clock();
			double cpu_time = (double(end_time - start_time)) / CLOCKS_PER_SEC;
			medtime += cpu_time;
			cout << "Процессорное время: " << cpu_time << endl;
			cout << endl;			
		}
		cout << "Среднее значение работы процессора: " << medtime / num;
	}
	return 0;
}

\end{lstlisting}

%2. Расстояние Дамерау-Левенштейна нерекурсивным методом используя матрицу.
\begin{lstlisting}[caption={Расстояние Дамерау-Левенштейна нерекурсивный метод, используя матрицу},label={lst:label2}]
#include <iostream>
#include <string>
#include <algorithm>
#include <vector>
#include <ctime>

using namespace std;

string RandomWord(int length) {
	static const char alphabet[] = "абвгдеёжзиклмнопрстуфхцчшщъыьэюя";
	string result;

	for (int i = 0; i < length; ++i) {
		int randomIndex = rand() % (sizeof(alphabet) - 1);
		result += alphabet[randomIndex];
	}

	return result;
}

int DamerauLevenshteinDistance(const string& source, const string& target) {

	int n = source.length();
	int m = target.length();
	vector<vector<int>> distanceMatrix(n + 1, vector<int>(m + 1));

	for (int i = 0; i <= n; ++i) {
		distanceMatrix[i][0] = i;
	}
	for (int j = 0; j <= m; ++j) {
		distanceMatrix[0][j] = j;
	}

	for (int i = 1; i <= n; ++i) {
		for (int j = 1; j <= m; ++j) {

			int substitutionCost = (source[i - 1] != target[j - 1]);

			int deletionCost = distanceMatrix[i - 1][j] + 1;
			int insertionCost = distanceMatrix[i][j - 1] + 1;
			int substitutionOrEqualityCost = distanceMatrix[i - 1][j - 1] + substitutionCost;

			//  транспозиция, если это возможно
			if (i > 1 && j > 1 && source[i - 1] == target[j - 2] && source[i - 2] == target[j - 1]) {
				distanceMatrix[i][j] = min({ deletionCost, insertionCost, substitutionOrEqualityCost, distanceMatrix[i - 2][j - 2] + substitutionCost });
			}
			else {
				distanceMatrix[i][j] = min({ deletionCost, insertionCost, substitutionOrEqualityCost });
			}
		}
	}

	for (int i = 0; i <= source.length(); ++i) {
		for (int j = 0; j <= target.length(); ++j) {
			cout << distanceMatrix[i][j] << " ";
		}
		cout << endl;
	}

	return distanceMatrix[n][m];
}

int main() {
	setlocale(0, "");
	string str1, str2;
	srand(time(0));
	int a, k, num;
	double medtime = 0;
	cout << "Ввести данные вручную (1) или сгенерировать (2) - ";
	cin >> a;
	if (a == 1) {
		cout << "Введите первое слово: ";
		cin >> str1;
		cout << "Введите второе слово: ";
		cin >> str2;
		clock_t start_time = clock();
		cout << "Расстояние Дамерау-Левенштейна: " << DamerauLevenshteinDistance(str1, str2) << endl;
		clock_t end_time = clock();
		double cpu_time = (double(end_time - start_time)) / CLOCKS_PER_SEC;
		cout << "Процессорное время: " << cpu_time << endl;
		cout << endl;
	}
	else {
		cout << "Введите длину слова: ";
		cin >> k;
		cout << "Количество слов: ";
		cin >> num;
		for (int i = 1; i <= num; i++) {
			cout << i << endl;
			str1 = RandomWord(k);
			str2 = RandomWord(k);
			cout << str1 << endl << str2 << endl;
			clock_t start_time = clock();
			cout << "Расстояние Дамерау-Левенштейна: " << DamerauLevenshteinDistance(str1, str2) << endl;
			clock_t end_time = clock();
			double cpu_time = (double(end_time - start_time)) / CLOCKS_PER_SEC;
			medtime += cpu_time;
			cout << "Процессорное время: " << cpu_time << endl;
			cout << endl;		
		}
		cout << "Среднее значение работы процессора: " << medtime / num;
	}
	return 0;
}
\end{lstlisting}

\begin{lstlisting}[caption={Расстояние Дамерау-Левенштейна рекурсивный метод, используя три строки матрицы},label={lst:label3}]
#include <iostream>
#include <string>
#include <vector>
#include <algorithm>
#include <ctime>

using namespace std;

string RandomWord(int length) {
	static const char alphabet[] = "абвгдеёжзиклмнопрстуфхцчшщъыьэюя";
	string result;

	for (int i = 0; i < length; ++i) {
		int randomIndex = rand() % (sizeof(alphabet) - 1);
		result += alphabet[randomIndex];
	}

	return result;
}

int damerauLevenshteinDistance(const string& str1, const string& str2, int i, int j, vector<vector<int>>& memo) {
	if (i == 0)
		return j;

	if (j == 0)
		return i;

	if (memo[i][j] != -1)
		return memo[i][j];

	int cost = (str1[i - 1] == str2[j - 1]) ? 0 : 1;

	int deletee = damerauLevenshteinDistance(str1, str2, i - 1, j, memo) + 1;
	int insert = damerauLevenshteinDistance(str1, str2, i, j - 1, memo) + 1;
	int substitution = damerauLevenshteinDistance(str1, str2, i - 1, j - 1, memo) + cost;

	int distance = min({ deletee, insert, substitution });

	if (i > 1 && j > 1 && str1[i - 1] == str2[j - 2] && str1[i - 2] == str2[j - 1]) {
		int transposition = damerauLevenshteinDistance(str1, str2, i - 2, j - 2, memo) + cost;
		distance = min(distance, transposition);
	}

	memo[i][j] = distance;
	return distance;
}

int damerauLevenshteinDistance(const string& str1, const string& str2) {
	int m = str1.length();
	int n = str2.length();

	vector<vector<int>> memo(m + 1, vector<int>(n + 1, -1));

	return damerauLevenshteinDistance(str1, str2, m, n, memo);
}

int main() {
	setlocale(0, "");
	string str1, str2;
	srand(time(0));
	int a, k, num;
	double medtime = 0;
	cout << "Ввести данные вручную (1) или сгенерировать (2) - ";
	cin >> a;
	if (a == 1) {
		cout << "Введите первое слово: ";
		cin >> str1;
		cout << "Введите второе слово: ";
		cin >> str2;
		clock_t start_time = clock();
		cout << "Расстояние Дамерау-Левенштейна: " << damerauLevenshteinDistance(str1, str2) << endl;
		clock_t end_time = clock();
		double cpu_time = (double(end_time - start_time)) / CLOCKS_PER_SEC;
		cout << "Процессорное время: " << cpu_time << endl;
		cout << endl;
	}
	else {
		cout << "Введите длину слова: ";
		cin >> k;
		cout << "Количество слов: ";
		cin >> num;
		for (int i = 1; i <= num; i++) {
			cout << i << endl;
			str1 = RandomWord(k);
			str2 = RandomWord(k);
			cout << str1 << endl << str2 << endl;
			clock_t start_time = clock();
			cout << "Расстояние Дамерау-Левенштейна: " << damerauLevenshteinDistance(str1, str2) << endl;
			clock_t end_time = clock();
			double cpu_time = (double(end_time - start_time)) / CLOCKS_PER_SEC;
			medtime += cpu_time;
			cout << "Процессорное время: " << cpu_time << endl;
			cout << endl;
		}
		cout << "Среднее значение работы процессора: " << medtime / num;
	}



	return 0;
}
\end{lstlisting}

\newpage
\begin{center}
    \section{Исследовательская часть}
\end{center}
\justifying

В таблице \ref{tab:my_label} приведено время в секундах, потраченное на вычисление расстояния определённым алгоритмом:
\begin{table}[h]
\begin{center}
    \centering
    \begin{tabular}{ | m{6em} | m{2cm}| m{2cm} | m{2cm} | m{2cm} | m{2cm} |} 
  \hline
 Алгоритм & 5 символов & 10 символов & 15 символов & 20 символов & 50 символов\\ 
  \hline
  Левенштейн без рекурсии 2 строки & 0.002 & 0.0026 & 0.0028 & 0.004 & 0.0097\\ 
  \hline
  Дамерау-Левенштейн без рекурсии матрица & 0.012 & 0.029 & 0.063 & 0.116 & 0.6396\\ 
  \hline
  Дамерау-Левенштейн рекурсия с кэшем & 0.002 & 0.003 & 0.004 & 0.0058 & 0.0252 \\
  \hline
\end{tabular}
\caption{\label{tab:my_label}}
\end{center}
\end{table}

Результаты вычислений также представленны на графике (рисунок ~\ref{img:grap1}): 
\begin{figure}[h!]
    \center{\includegraphics{graph}}
    \caption{График зависимости затраченного времени}
    \label{img:grap1}
\end{figure}

\justifying
Нетрудно заметить, что затраченное время возрастсает с увеличением количества вычислений. Поскольку при заполнении матрицы требуется сделать $3^{n \cdot n}$ вычислений, где n - количество символов в слове, то необходимо расходовать большое количество времени. Что касается алгоритма Дамерау-Левенштейна с рекурсией с кэшем, то затраченного времени и расходованной памяти, безусловно, становится гораздо меньше. Однако при больших объемах информации лидирует алгоритм Левенштейна без рекурсии с помощью двух строк. Данный алгоритм универсален, так как хранит в разы меньше информации по сравнению с другими алгоритмами и тратит на вычисления меньше времени. Он будет являться самым ``выигрышным''.

\newpage
\begin{center}
    \section*{Заключение}
    \addcontentsline{toc}{section}{Заключение}
\end{center}

%В результате выполнения лабораторной работы были реализованы такие алгоритмы, как расстояние Левенштейна без рекурсии в две строки, расстояние Дамерау-Левенштейна без рекурсии, используя матрицу, расстояние Дамерау-Левенштейна с рекурсией с кэшем. Для удобства нахождения эффективного алгоритма был создан автоматический генератор слов разной длины. После обработки полученных данных наиболее эфективным алгоритмом оказался Левенштейна без рекурсии в две строки. Он меньше всего затрачивает память компьютера, что делает его максимально эффективным по сравнению с другими алгоритмами.
В результате выполнения лабораторной работы:
\begin{enumerate}
\item Описано расстояние Левенштейна и Дамерау-Левенштейна.
\item Разработаны алгоритмы вычисления редакционных расстояний, согласно варианту.
\item Реализованы разработанные алгоритмы.
\item Выполнено тестирование реализации разработанных алгоритмов.
\item Разработано программное обеспечение с двумя режимами работы --- режимом расчёта расстояний для введённых пользователем двух строк произвольной длины и режимом массированного замера процессорного времени для автоматически сгенерированных строк заданной длины.
\item Проведено сравнение временной эффективности реализации разработанных алгоритмов.
\item Выполнена оценка емкостной эффективности разработанных алгоритмов в зависимости от линейных размеров строк.
\end{enumerate}

Согласно варианту, разработанно следующие три алгоритма.
\begin{enumerate}
\item Алгоритм вычисления расстояния Левенштейна нерекурсивным методом, используя две строки.
\item Алгоритм вычисления расстояния Дамерау-Левенштейна нерекурсивным методом, пользуясь матрицей и в конце программы вывести её содержимое.
\item Алгоритм вычисления расстояния Дамерау-Левенштейна рекурсивным методом без кэша.
\end{enumerate}
После обработки полученных данных наиболее эфективным алгоритмом оказался Левенштейна без рекурсии в две строки. Он меньше всего затрачивает память компьютера, что делает его максимально эффективным по сравнению с другими алгоритмами.

\newpage
\begin{center}
\begin{thebibliography}{}
\addcontentsline{toc}{section}{Список литературы}
\bibitem{mpl} Левенштейн В. И. Двоичные коды с исправлением выпадений, вставок
и замещений символов / Доклады АН СССР, т. 163. — М.: Издательство «Наука», 1965.
\end{thebibliography}
\end{center}
\end{document}
